\documentclass[a4paper,twoside]{article}

\usepackage{morewrites}

\usepackage[svgnames,dvipsnames,x11names]{xcolor}
\usepackage{graphicx}
\usepackage[english]{babel}
\usepackage[colorlinks=true, linkcolor=NavyBlue, urlcolor=NavyBlue, citecolor=NavyBlue]{hyperref}
\usepackage[a4paper]{geometry}
\usepackage{ts-typesetting}
\usepackage{ts-fonts}
\usepackage{csquotes}
\usepackage{xcolor}
\usepackage{epigraph}
\usepackage{marginnote}
\usepackage{multicol}
\usepackage{progressbar}
\usepackage{graphicx}
\usepackage{float}
\usepackage{seqsplit}
\renewcommand*{\marginfont}{
\footnotesize
\sloppy
\emergencystretch=1em
\hyphenpenalty=50
\exhyphenpenalty=50
}
\newlength{\tsmarginparavail}
\newlength{\tsmarginparalt}
\newlength{\tsmarginparbuf}
\setlength{\tsmarginparbuf}{6mm}
\makeatletter
\AtBeginDocument{
\setlength{\tsmarginparavail}{\dimexpr\paperwidth-\textwidth-1in-\oddsidemargin-\marginparsep-\tsmarginparbuf\relax}
\ifdim\tsmarginparavail<\z@\setlength{\tsmarginparavail}{\z@}\fi
\if@twoside
\setlength{\tsmarginparalt}{\dimexpr1in+\evensidemargin-\marginparsep-\tsmarginparbuf\relax}
\ifdim\tsmarginparalt<\z@\setlength{\tsmarginparalt}{\z@}\fi
\ifdim\tsmarginparalt<\tsmarginparavail
\setlength{\tsmarginparavail}{\tsmarginparalt}
\fi
\fi
\ifdim\tsmarginparavail<\marginparwidth
\setlength{\marginparwidth}{\tsmarginparavail}
\fi
}
\makeatother

\usepackage{lastpage}
\usepackage{refcount}
\makeatletter
\def\ps@plain{
\let\@oddhead\@empty
\let\@evenhead\@empty
\def\@oddfoot{
\hfil
\ifnum\getpagerefnumber{LastPage}>1\relax
\thepage
\fi
\hfil}
\let\@evenfoot\@oddfoot
}
\makeatother

\title{About}
\author{}\date{}
\begin{document}
\maketitle
TeXSmith was originally created by Yves Chevallier in 2025 to address the need for a seamless workflow between Markdown-based documentation and \LaTeX{}-based publishing, initially for his own academic \href{https://heig-tin-info.github.io/handbook/}{courses} at HEIG-VD.

Aside from \href{https://pandoc.org/}{Pandoc}---written in \href{https://en.wikipedia.org/wiki/Haskell}{Haskell} and not directly suited for \href{https://www.mkdocs.org/}{MkDocs}---there were no tools capable of converting MkDocs-flavored Markdown into \LaTeX{} while preserving the original content's semantic intent.

Because developing such an ambitious toolchain was a substantial and time-consuming effort, I postponed the project until I discovered the remarkable power of OpenAI Codex, which helped me bootstrap the initial version of TeXSmith in just a few days. I wanted to extract the core MkDocs-to-\LaTeX{} code used in my online course and turn it into a standalone, general-purpose tool that anyone needing to convert MkDocs content to \LaTeX{} could use. That is how TeXSmith was born.
\section{Branding}
This project is \textbf{not affiliated with \TeX{}, \LaTeX{}, or the \LaTeX{} Project}. It merely produces \LaTeX{}-compatible output and interacts with the \TeX{} toolchain in the same way any document-generation utility would. All trademarks belong to their respective owners.

By convention within the \TeX{} community, the names \emph{\TeX{}} and \emph{\LaTeX{}} are used with care. \href{https://en.wikipedia.org/wiki/Donald\_Knuth}{Donald Knuth} famously stated:
\begin{displayquote}
\tslead{``\TeX{} is not to be changed; only Knuth himself may change \TeX{}.''}
\end{displayquote}
This is not a legal trademark declaration but a long-standing cultural rule: any system that calls itself \emph{\TeX{}} must be fully compatible with Knuth's canonical implementation. Similarly, the \LaTeX{} Project requires that only implementations conforming to the \LaTeX{} format may use the name \emph{\LaTeX{}}.

In keeping with these established norms, this project does \textbf{not} claim to be a \TeX{} or \LaTeX{} implementation, nor does it modify or replace them. It is simply a tool that \emph{generates} \LaTeX{} code as output, leaving the actual typesetting to standard, community-maintained engines.
\end{document}
